\scp{\tsc{Nuclear \ili{Central} Timor}}{05-03}
Within \tsc{\ili{Central} Timor} there is evidence for a \tsc{Nuclear
\ili{Central} Timor} subgroup which contains all languages except \ili{Welaun}.
Evidence for this subgroup comes from lowering of high vowels
to mid before  final *R,
the merger of **mb/ **mp > \it{b} in most cases (as discussed
above), and penultimate *ə > \it{e}.
Firstly, \tsc{Nuclear \ili{Central} Timor} languages usually have lowering
of high vowels to mid before final *R,
which was then often lost.
Examples are given in \trf{05-11}.

\begin{table}\tcp{\tsc{\ili{Central} Timor} *-uR{\#} and *-iR{\#}}{05-11}
\begin{adjustbox}{width=\textwidth}
	\begin{threeparttable}\st{2pt}
	\begin{tabular}{llllll}\lsptoprule
local gl.	&	tail	&	egg	&	grain head	&	dust, ashes\su{Δ}	&	breadfruit\\
PMP	&	*ik\tbr{uR}	&	*qatəl\tbr{uR}	&	*bul\tbr{iR}	&	*dap\tbr{uR}	&	*kul\tbr{uR}\\ \midrule
\ili{Welaun}	&	{\hr}hiʔ\tbb{u}-n\su{†}	&	{\hr}man tol\tbb{u}-n	&	{\hr}ful\tbb{i}n	&	\cg{(lai lahuk-aan)}	&	{\hr}kul\tbb{u}\\
\ili{Kemak}	&	{\hr}hiʔ\tbr{o}-r\su{†}	&	{\hr}manu tel\tbr{o}-n	&	{\hr}hul\tbr{e}-n	&	{\hr}rap\tbr{o}, rap\tbr{o} apan	&	{\hr}kul\tbr{o}\\
\ili{Tokodede}	&	{\hr}ik\tbr{o}	&	{\hr}manu tel\tbr{o}-n	&		&	{\hr}rap\tbr{o}	&	{\hr}kul\tbr{o}r\\
NW. \ili{Mambae}	&	{\hr}i\tbr{o}-n	&	{\hr}man tel\tbr{o}-n	&		&		&	\\
C. \ili{Mambae}	&	{\hr}i\tbr{o}-n	&	{\hr}maun tel\tbr{o}-n	&		&		&	\\
S. \ili{Mambae}\su{†}	&	{\hr}i\tbr{o}	&	{\hr}man tel\tbr{o}	&		&	{\hr}aap rap\tbr{o}	&	\\
		\lspbottomrule
	\end{tabular}
		\begin{tablenotes}\tnsize
			\item[†]	{South \ili{Mambae} data from Hatu-Udo subdistrict.}
			\item[‡]	{\ili{Welaun} \it{hiʔu-n}
								and \ili{Tokodede} \it{hiʔo-r} ‘tail’ have insertion of /h/ /{\#\gap}VʔV.}
			\item[Δ]	{\ili{Kemak} \it{rapo} is ‘dust; slaked lime,
								calcium’ and \it{rapo apan} is ‘ash in the fireplace’.
								\ili{Tokodede} \it{rapo} is only given as ‘lime, calcium’ by \citet[23]{LSG2006}.
								South \ili{Mambae} \it{aap rapo} means ‘ashes’.
								(Initial \it{aap} ‘fire’ is metathesised from **api.) South \ili{Mambae}
								from Letefoho village has \it{aep ni rapu} with final *u = \it{u} /{\gap}R{\#}.
								The semantic shift from PMP *dapuR ‘hearth’ > ‘calcium’ seems
								motivated by the similarity between the ash in the hearth
								and powdered calcium, as is used for chewing betel nut.}
		\end{tablenotes}
	\end{threeparttable}
\end{adjustbox}
\end{table}

Secondly, all \tsc{\ili{Central} Timor} languages except \ili{Welaun} have
penultimate *ə > \it{e}.
\ili{Welaun} has penultimate *ə > \it{o} in most cases,
which is also the reflex in \ili{Tetun} (\srf{03-03}).
(One exceptional word in \ili{Welaun} is *ənəm > \it{inam} ‘six’.)
Examples of *ə are given in \trf{05-12}.

\begin{table}\tcp{\tsc{\ili{Central} Timor} non-final *ə}{05-12}
\begin{adjustbox}{width=\textwidth}
	\begin{threeparttable}\st{2pt}
	\begin{tabular}{llllllll}\lsptoprule
local gl.	&	three	&	corn\su{†}	&	sun	&	new	&	egg	&	hear	&	flea\\
PMP	&	*t\tbr{ə}lu	&	*z\tbr{ə}lay	&	*qal\tbr{ə}jaw	&	*baq\tbr{ə}Ru	&	*qat\tbr{ə}luR	&	*d\tbr{ə}ŋəR	&	*qatim\tbr{ə}la\\ \midrule
\ili{Welaun}	&	{\hr}t\tbr{o}lu	&	{\hr}s\tbr{o}le	&	{\hr}l\tbr{o}ro	&	{\hr}f\tbr{o}unaan	&	{\hr}man t\tbr{o}lu-n	&	{\hr}l\tbr{o}ka	&	{\hr}tam\tbr{o}la\\
\ili{Kemak}	&	{\hr}t\tbr{e}lu	&	{\hr}s\tbr{e}le	&	{\hr}l\tbr{e}lo	&	{\hr}h\tbr{e}un	&	{\hr}manu t\tbr{e}lo-n	&	{\hr}r\tbr{e}ga	&	{\hr}tə̆m\tbr{e}la\\
\ili{Tokodede}	&	{\hr}t\tbr{e}lu	&	{\hr}s\tbr{e}le	&	{\hr}l\tbr{e}lo	&	{\hr}h\tbr{e}uː	&	{\hr}manu t\tbr{e}lo-n	&	\cg{(pliʔi)}	&	{\hr}tum\tbr{e}l\\
NW. \ili{Mambae}	&	{\hr}t\tbr{e}ul	&	{\hr}s\tbr{e}la	&	{\hr}l\tbr{e}lon	&	{\hr}h\tbr{e}u	&	{\hr}man t\tbr{e}lo-n	&	\cg{(pliik)}	&	\\
C. \ili{Mambae}	&	{\hr}t\tbr{e}ul	&	{\hr}s\tbr{e}la	&	{\hr}l\tbr{e}el-a	&	{\hr}h\tbr{e}un	&	{\hr}maun t\tbr{e}lo-n	&	\cg{(fliik)}	&	\\
S. \ili{Mambae}	&	{\hr}t\tbr{e}ul	&	\cg{(batar)}	&	{\hr}l\tbr{e}ol mata	&	{\hr}h\tbr{e}u	&	{\hr}maun t\tbr{e}lo	&	\cg{(fliik)}	&	\\
		\lspbottomrule
	\end{tabular}
		\begin{tablenotes}\tnsize
			\item[†] {PMP *zəlay is reconstructed with the meaning ‘Job’s tears’.}
		\end{tablenotes}
	\end{threeparttable}
\end{adjustbox}
\end{table}

\tsc{Nuclear \ili{Central} Timor} languages also share medial *j >
\it{l} in five words, given in \trf{05-13}.
While \ili{Welaun} also has *j > \it{l} in three of these words,
it has *j > \it{r} adjacent to schwa.
This, combined with the change of \ili{Welaun} *d > **r > \it{l},
indicates that \ili{Welaun} *j > \it{l} in these words may have gone
though intermediate **d and/or **r,
in which case the change of *j > \it{l} would not be shared between
all \tsc{\ili{Central} Timor} languages.\footnote{
		Note, however, that \ili{Welaun} has *d > \it{l} adjacent to schwa as shown
		by *dəŋəR > \it{loka} ‘hear’.
		The only other known cases of *d adjacent to schwa are *dəpa
		> \it{roʔa} ‘stretch out the arm,
		arm-span’, which may be borrowed from West \ili{Tetun} \it{roʔa},
		and *sumaŋəd > \it{makar-aat} ‘soul’,
		for which *d is word final.}

\begin{table}\tcp{\tsc{Nuclear \ili{Central} Timor} *j > \it{l}}{05-13}
%\begin{adjustbox}{width=\textwidth}
	\begin{threeparttable} %\st{2pt}
	\begin{tabular}{llllll}\lsptoprule
local gl.	&	ySib\su{†}	&	nose	&	how much	&	name	&	bitter\\
PMP	&	*hua\tbr{j}i	&	*i\tbr{j}uŋ	&	*pi\tbr{j}a	&	*ŋa\tbr{j}an	&	**ma-pə\tbr{j}u\su{Δ}\\	\midrule
\ili{Welaun}	&	{\hr}wa\tbr{l}i-n	&	{\hr}i\tbr{l}uk-aat	&	{\hr}hi\tbr{r}a	&	{\hr}ka\tbr{l}an-aat	&	{\hr\hr}fo\tbr{r}u\\
\ili{Kemak}	&	\hp{*w}a\tbr{l}i-r	&	{\hr}i\tbr{l}u-r	&	{\hr}pi\tbr{l}a	&	{\hr}ga\tbr{l}a-r	&	{\hr\hr}pe\tbr{l}un\\
\ili{Tokodede}	&	\hp{*w}a\tbr{l}i	&	{\hr}i\tbr{l}u	&	{\hr}pii\tbr{l}	&	{\hr}ga\tbr{l}aː	&	{\hr\hr}pe\tbr{l}u\\
NW. \ili{Mambae}	&	\hp{*w}a\tbr{l}i-n	&	{\hr}i\tbr{l}u-n	&	{\hr}arpii\tbr{l}\su{‡}	&		&	\\
C. \ili{Mambae}	&	\hp{*w}a\tbr{l}i-n	&	{\hr}i\tbr{l}u-n	&	{\hr}arfii\tbr{l}\su{‡}	&		&	\\
S. \ili{Mambae}	&	\hp{w}\cg{(kau)}	&	{\hr}i\tbr{l}u	&	{\hr}arpii\tbr{l}\su{‡}	&	{\hr}ka\tbr{l}a	&	\\
		\lspbottomrule
	\end{tabular}
		\begin{tablenotes}\tnsize
			\item[†]	{Reflexes of *huaji mean ‘same sex younger sibling’ in \tsc{\ili{Central} Timor} languages.}
			\item[‡]	{The \ili{Mambae} words \it{arpiil, arfiil} mean ‘when?’.}
			\item[Δ]	{**ma-pəju ‘bitter’ is from PMP *qapəju ‘gall, gallbladder, bile’.}
		\end{tablenotes}
	\end{threeparttable}
%\end{adjustbox}
\end{table}

Apart from the five words which have *j > \it{l} in \tsc{Nuclear
\ili{Central} Timor}, the other known reflexes of *j are: PMP *ma\tbr{j}a ‘fry’ > \ili{Welaun}
\it{ma\tbr{r}a}; \ili{Kemak} \it{ma\tbr{l}an}; South \ili{Mambae} **masa
> \it{maa\tbr{s}}, *bu\tbr{j}əq ‘foam’ > \ili{Welaun} \it{fo\tbr{r}en}; \ili{Kemak} \it{bia hu\tbr{l}an}; \ili{Tokodede} \it{hu\tbr{r}iː},
*pusə\tbr{j} ‘navel; centre’ > \ili{Welaun} \it{husa\tbr{r}-aat};
\ili{Atabae} \ili{Kemak} \it{pus\tbr{r}-aar}; \ili{Tokodede} \it{kupusa};
South \ili{Mambae} \it{fusu}, and probably *qulə\tbr{j} > \ili{Kemak} \it{ula\tbr{l}}, \ili{Tokodede}
\it{ula} ‘snake’, though this last word may be from *hulaR.
